% !TEX TS-program = pdflatexmk
\documentclass[12pt]{article}

% Layout.
\usepackage[top=1in, bottom=0.75in, left=1in, right=1in, headheight=1in, headsep=6pt]{geometry}

% Fonts.
\usepackage{mathptmx}
\usepackage[scaled=0.86]{helvet}
\renewcommand{\emph}[1]{\textsf{\textbf{#1}}}

% Misc packages.
\usepackage{amsmath,amssymb,latexsym}
\usepackage{graphicx,hyperref}
\usepackage{array}
\usepackage{xcolor}
\usepackage{multicol}
\usepackage{tabularx,colortbl}
\usepackage{enumitem}

\hypersetup{
    colorlinks=true,
    linkcolor=blue,
    filecolor=magenta,      
    urlcolor=blue,
    pdftitle={Calc I Syllabus},
    pdfpagemode=FullScreen,
    }

\def\mailto#1{\href{mailto:#1}{#1}}

% Paragraph spacing
\parindent 0pt
\parskip 6pt plus 1pt
\def\tableindent{\hskip 0.5 in}
\def\ts{\hskip 1.5 em}

\usepackage{fancyhdr}
\pagestyle{fancy} 
\lhead{\large\sf\textbf{MATH F251X: Calculus I (in person)}}
\rhead{\large\sf\textbf{Spring 2026 Syllabus}}

\newcommand{\localhead}[1]{\par\smallskip\textbf{#1}\nobreak\\}%
\def\heading#1{\localhead{\large\emph{#1}}}
\def\subheading#1{\localhead{\emph{#1}}}

\newenvironment{clist}%
{\bgroup\parskip 0pt\begin{list}{$\bullet$}{\partopsep 4pt\topsep 0pt\itemsep -2pt}}%
{\end{list}\egroup}%

%%Commenting tools. You can ignore this stuff unless we're exchanging materials where I'm commenting in detail on how you are writing/using LaTeX.
\usepackage{soul,cancel,color}
\usepackage[normalem]{ulem}
\usepackage[dvipsnames,usenames]{xcolor}
\newcommand{\red}[1]{{\color{BrickRed}#1}}
\newcommand{\blue}[1]{{\color{MidnightBlue}#1}}
\newcommand{\green}[1]{{\color{OliveGreen}#1}}
\newcommand{\medit}[2]{\red{\cancel{#1}}{\colorbox{yellow!100}{\green{#2}}}} 
\newcommand{\edit}[2]{\red{\sout{#1}}\green{\hl{#2}}}
\definecolor{OliveGreen}{rgb}{.3,.5,.2}
\definecolor{MidnightBlue}{rgb}{.3,.4,.6}
\newcommand{\del}[1]{\red{\sout{#1}}}
\newcommand{\inst}[1]{\green{\hl{#1}}}
\newcommand{\minst}[1]{\colorbox{yellow!100}{$\green{#1}$}}
\usepackage[displaymath, mathlines]{lineno}




\begin{document}

% \heading{Course Description}

\strut\par\vskip-12pt
\heading{Essential Information}

\vskip -12pt
\strut\hbox to \hsize{\tableindent\vtop{\halign{#\hfill\ts&#\hfil\cr
{\emph{Instructors}}&Dr. James Gossell (\S 901 \& \S 001), \& Dr. Gordon Williams (\S 902), \cr
\strut & \cr

{\emph{Website}}&\url{https://uaf-math.github.io/calc1/}\cr
\strut & \cr
\emph{Prerequisite} & MATH F151X and MATH F152X; or MATH F156X; or placement.\cr
\strut & \cr
{\emph{Required Text}} &\textit{OpenStax Calculus Volume 1} by G. Strang \& E. Herman,\cr
&\url{https://openstax.org/details/books/calculus-volume-1} \cr
&(optional print copy)\texttt{ISBN-13: 978-1938168024}\cr
  \strut & \cr
{\emph{Grades}}&(Canvas) \url{https://www.uaf.edu/uaf/current/canvas.php}\cr
 \strut & \cr
%{\emph{Class Discord}}&\url{https://nookbot.uaf.edu/}\cr
 \strut & \cr
 

  {\emph{Dr. James Gossell}}&Email: jegossell@alaska.edu\cr
 &Office: Chapman 301D\cr
  \strut & \cr 

  {\emph{Dr. Gordon Williams}}&Email: giwilliams@alaska.edu\cr
 &Office: Chapman 306D\cr
  \strut & \cr
 }
\hfil}}

\vskip -12pt
\heading{Description, Course Goals \& Student Learning Outcomes}
Calculus is one of mathematics' premiere computational tools.  It has
pervasive applications in all the sciences and is part of the UAF core 
curriculum.  The two principal tools of calculus are differentiation and integration. Differentiation concerns how changes in one variable affect another.
How does a population of bacteria change as time changes?  
How does the temperature of the ocean change as depth increases?
Integration, on the other hand, is a kind of reverse process to
differentiation. 

Students completing the course will have the mathematical foundation to be successful in Calculus II and other courses
requiring this background.  Specifically, students will
\begin{clist}
\item {demonstrate an} understand{ing of} the role of limits in the definition of a derivative{, and be able to evaluate a derivative in this way} 
and be able to compute elementary derivatives from this definition,
\item {demonstrate }understand{ing of} the definition of a continuous function {by} identify{ing}
continuous/discontinuous functions,
\item compute standard derivatives,
\item apply derivatives to common types of applied problems,
\item {explain} the definition of the the definite integral,
\item apply the Fundamental Theorem of Calculus to
compute definite integrals,
\item apply integration to common types of applied problems.
\end{clist}

\newpage

\heading{Class Time}
There are four hours of class meetings with your primary instructor every week (MWThF). This time will be used to discuss new topics in Calculus and for a weekly quiz on Thursdays. Starting in Week 2, the last 30 minutes of the Thursday class will be used to go over the quiz so all students can leave knowing what material on that quiz has been mastered, how to work each problem correctly,  and what topics need additional study. 

The Tuesday Recitation (MATH F251L) will initially be spent helping students get started on first-of-semester tasks. Starting on Week 3, the Tuesday class is explicitly devoted to bolstering the underlying non-Calculus skills that are nevertheless essential to success in Calculus such as: graphing, algebra, trigonometry, exponential and logarithmic functions, and inverse functions. It will also include additional strategic homework/quiz/test prep. As a concrete example, one of the things we will do in Recitation is go over the non-Calculus math skills needed to complete the up-coming homework so that students can focus on the Calculus instead of getting bogged down in algebra or trigonometry.

\heading{Tentative Schedule}
The course website contains a %\href{https://docs.google.com/spreadsheets/d/e/2PACX-1vQGc6J1T0q0ik53Vks8LatvLNG9OMWIlA9mkIOz6BWO9KZecpv0yYPTKDhgKxS-UxQ809bjz4kUVgbk/pubhtml?gid=0&single=true}{schedule}
schedule for the semester listing
the topics to be covered each class, the dates each assignment is due,
the topics of every quiz, and so forth. You should consult this schedule
routinely.  {\bf We may make minor adjustments to the schedule, which
will be announced in advance.}

\heading{Office Hours and Communication}
Instructors and TAs will schedule formal office hours, which will be listed on web sites accessible from the main course webpage. Any office hours in the \href{https://uaf.edu/dms/mathlab/}{Math Tutoring Lab} are open to \emph{all} Calculus I students regardless of who is their instructor of record. Students can also schedule meetings with their instructor outside of regular office hours.

We will use Canvas to send announcements. If we (your instructor/TA) need to contact you, we will first try to do this in class. Our second method will be to send an email to you via Canvas. Thus, you will want to make sure that the email address in Canvas is one that you check regularly. Note that in Canvas it is possible to set up text alerts. However, you must login to Canvas and adjust the setting for your account. Neither email nor text alerts are automatic.



\heading{Evaluation and Grades}
Grades are determined as follows; 
each component of the grade is discussed subsequently in
the syllabus.
 
\begin{multicols}{2}

\begin{tabular}{|c|c|}
\hline
Class Participation& 7.5\%\\
\hline
(Written) Homework& 7.5\% \\
\hline
Quizzes& 15\% \\
\hline
Midterm 1 & 15\% \\
\hline
Derivative Proficiency& 7.5\%\\
\hline
Midterm 2 & 15\%  \\
\hline
Integral Proficiency& 7.5\%\\
\hline
Final Exam& 25\% \\
\hline
Total& 100\%\\
\hline
\end{tabular}

\vskip 6pt
Letter grades will be assigned according to the following scale.
This scale is a guarantee; the instructors reserve the right to lower the thresholds\footnote{{This would improve your grade.}}. 

\def\sts{\hskip 0.5em}
\strut\hbox to\hsize{\vbox{\halign{#\hfil\sts&#\hfil\ts&#\hfil\sts&#
\hfil\ts&#\hfil\sts&#\hfil   \cr
A+ & 97--100\% & C+ & 77--79.99\% & F  & $<$ 60\%\cr

A & 93--96.99\% &  C & 70--76.99\%&&\cr
A- & 90--92.99\% & C- & not given&&\cr
B+ & 87--89.99\% & D+ & 67--69.99\%&&\cr
B &  83--86.99\% & D & 63--66.99\%&&\cr
B- & 80-82.99\% & D- & 60--62.99\%&&\cr
}}\hfil}
\end{multicols}

\pagebreak

\heading{Online Course Materials}
All materials for this course are posted online. (See below.) \\

\begin{table}[h]
\centering
\begin{tabular}{rl}
Where to find it&What you are looking for\\
\hline \hline
&syllabus and day-by-day schedule for the semester\\
&ALEKS logistics\\
&homework problem sets\\
public Calculus I webpage&practice quizzes, proficiencies, midterms, final exams\\
\href{https://uaf-math.github.io/calc1/}{https://uaf-math.github.io/calc1/}&solutions to practice quizzes, proficiencies, midterms, final exam\\
&videos\\
&recitation worksheets with solutions \& videos\\
&textbook\\
&in-class worksheets with solutions\\
\hline
&announcements and reminders\\
&link to Gradescope (to turn in homework)\\
\href{https://www.uaf.edu/uaf/current/canvas.php}{Canvas}& complete solutions to homework\\
& your grades\\
&textbook\\
%\hline
%see your instructor&in-class materials for your section\\
\end{tabular}
%\caption{}
%\label{wheretofind}
\end{table}

%\heading{Weeks 1 \& 2 Logistics}
%The first week of the course is devoted to prerequisite review. Consequently, the homework and quiz mechanics for the first two weeks are different from that of the remainder of the semester.
%
%In the first week, instead of the usual written homework, you will be
%working with a program called ALEKS to refresh precalculus skills, and the first
%quiz will be an ALEKS-based assessment. (See \href{https://uaf-math251.github.io/week1-ALEKS.html}{Week 1 Details} for step-by-step instructions.)\\ 
%During the first weeks you will:
%\begin{clist}
%\item enroll in the UAF MATH 251 Spring 2022 Cohort of ALEKS PPL\\
%class code: \textbf{ERD3E-L6JKM}
%\item complete an initial placement test (approx 1-2 hours) by \textbf{Tuesday August 30} at 11:59pm \\
%(Note: We will start this in class on Tuesday. Most people will finish during class. Those who do not will just finish up later that day. Finishing the initial placement on time is worth 10 points added to Quiz 1, your proctored ALEKS score.),
%\item complete 90\% of the ALEKS pie \: OR \: spend 5 hours in Learning Mode by Monday, September 5 at 11:59pm which will count as your first homework grade,
%\item sign-up by Monday, September 5 for a 2-hour Tuesday time-slot to take Quiz 1, 
%\item complete an ALEKS assessment (approx 1-2 hours) on Tuesday, September 6 which will count as Quiz 1.
%\end{clist}
%The standard Tuesday class on September 6 is cancelled and instead you will go to Rasmuson Library Room 301 to take your ALEKS assessment at whatever time slot you reserved. 

\heading{Class Participation}
Attendance and participation in class is an important part of mastering the material in Calculus (and all of your classes). Class attendance is one of the best predictors of overall course performance regardless of subject. For this reason, we want to incentivize this aspect of your education. A student is counted as having successfully participated in class if that student was in class on-time and was an active participant during the whole class. Being an active participant includes engaging in classroom tasks such as individual or group worksheets, classroom discussions, and other classroom activities. Your class participation score will be calculated as the percent of classes in which you were an active participant. Let your instructor know if you have {or will} missed class for an excus{able} reason.

\heading{Homework}
Homework assignments typically consist of a selection of problems at
the end of each section of our textbook and a few additional
problems. Homework is written (on paper or tablet) and turned in via
Gradescope which is accessed from Canvas.  Help with scanning homework
can be found under
\href{https://uaf-math.github.io/calc1/techHelp.html}{Technology Help} on
the course webpage. Assignments are due most Mondays and Wednesdays
(by 11:59 PM) in advance of the Thursday quiz. Complete worked
solutions to all problems are provided in advance on
Canvas. Consequently, your homework will be graded based on
\emph{effort} and \emph{completion}. Homework can be turned in up to
three days late with no penalty but will not be accepted after that.
It is your responsibility to seek help and complete homework
assignments on time.
 All students should earn
100\% of their homework points!

The list of homework problems and homework guidelines can be found at the \href{https://uaf-math.github.io/calc1/newhomework.html}{Homework} link on the course webpage. Your first regular written homework assignment is due \emph{Thursday, January 15.}

Clearly, it is possible to short-circuit the homework by copying the solutions. It should also be clear that (a) this is a bad idea and (b) your instructor/TA will know you have done this. Our goal in providing answers and solutions is to foster the use of homework as a \emph{learning experience}. 



\heading{Quizzes}
Quizzes 1 through 10 are paper quizzes and are given on Thursdays in the middle third of the class. Starting in Week 2, students will be given the opportunity to correct their quizzes in the last third of the Thursday class and can earn points back on their quiz for doing so \emph{accurately}. The weekly quiz will cover the material taught in the classes held since the previous quiz; specific topics can be found in the schedule on the course website. Additionally, there will be a proctored ALEKS placement given during recitation in Week 2 which tests Calculus readiness.

Quizzes are given under testing conditions; books, notes, and calculators are not allowed unless otherwise stated. Your performance on the quizzes \emph{prior to corrections} is the best indicator of how well you are learning the course material and much more accurate than your homework score.

The ALEKS quiz is a special quiz that is given during the second week of class. You will sign-up for a 2-hour time window to take your assessment. If you like, you can login to your ALEKS course in advance and/or see what the process will be like by visiting the ALEKS info link on the public webpage. The ALEKS quiz score will be scaled to be out of 25 points. All other quizzes will be worth 25 points. Students who take the ALEKS quiz will have their lowest quiz score dropped at the end of the semester.

\heading{Recitations}
The Tuesday Recitation is taught by a teaching assistant (TA), graduate students from the Math  Department who have experience teaching Calculus. Starting in Week 3, the Tuesday class is explicitly devoted to bolstering the underlying non-Calculus skills that are nevertheless essential to success in Calculus such as: graphing, algebra, trigonometry, exponential and logarithmic functions, and inverse functions. It will also include additional strategic homework/quiz/test prep. As a concrete example, one of the things we will do in Recitation is go over the non-Calculus math skills needed to complete the up-coming homework so that students can focus on the Calculus instead of getting bogged down in algebra or trigonometry.

Students may opt out of the recitations if they score 80 or higher on the Tuesday, {January 20} ALEKS assessment \textbf{and} if they have no unexcused absences during the first week. If such a student chooses to opt out of recitations, they are still allowed and encouraged to attend them, but a recitation absence will not count against them for their Class Participation portion grade. To opt out, the student must contact their instructor after Week 2 of the semester.

\heading{Midterms}
There are two midterm exams this semester, to be held on Thursday, February 12 and Thursday, April 9. Note that the course webpage contains all previous Midterms (with solutions) to provide lots of opportunity for students to practice. The midterms are the same for all sections; they are prepared and approved by all instructors teaching the course. 

\heading{Proficiencies}
A proficiency is an assessment covering a routine mechanical skill.  In this course we will take a Derivative Proficiency (March 5) and an Integral Proficiency (April 23). Note that the course webpage contains all previous proficiencies (with solutions) so a student can know in advance what these are like and has lots of opportunity for practice. Proficiencies will be graded on a binary scale for each problem (no partial credit). We have designed the grading structure in this course to prioritize and reward effort. Thus, you will have the opportunity to retake each proficiency. The highest of the scores will be used to calculate a student's grade.

%\pagebreak

\heading{Final Exam} 
The cumulative final exam will be held on Wednesday April 29. Note that the course webpage contains all previous final exams (with solutions) so a student can know in advance what these are like  and has lots of opportunity for practice. A make-up final exam will be given only in extenuating circumstances, for documented and excused reasons at the discretion of the instructors.

\heading{Make-ups}
Make-up quizzes and exams may be provided at the instructor's discretion only if
the student shows documentation supporting an excused absence. Excused
absences include absences due to illness, travel on University or
career-related business, and cultural or religious
observations. Unless it is an unforeseen circumstance, arrangements
must be made in advance; if unforeseen, the instructor must be notified
by the following day.

\heading{Tutoring and Resources}
\vskip -30pt\strut
\begin{clist}
	\item The Student Success Center on the 6th floor of the Rasmuson Library offers free drop-in tutoring.
	\item One-on-one (or small group) tutoring is available in 
Chapman Building Room 201. You must schedule an
appointment; see \url{https://www.uaf.edu/dms/mathlab/}.
	\item Student Support Services offers free tutoring in many subjects to students who qualify for their program.
	\item ASUAF offers private tutoring for a small fee (based on student income).
\end{clist}

\heading{Rules and Policies}
\vskip -20pt
%\subheading{Participation and Attendance}
%Class and recitation attendance is mandatory. Students who stop participating in the course will be withdrawn. Examples of inadequate participation include,
%but are not limited to:
%\begin{clist}
%\item missing class five times
%\item not completing or not turning in \textbf{three} written homework assignments
%\item failing to participate in classroom activities
%\item repeatedly failing tests and quizzes with no attempt at remediation
%\end{clist}

%\subheading{Disability Services}
%The Office of Disability Services implements the
%Americans with Disabilities Act (ADA), and ensures that UAF students
%have equal access to the campus and course materials. The instructors will work with
%the Office of Disability Services (208 Whitaker, 474-5655) to provide
%reasonable accommodations to students with disabilities.

%\subheading{Student Protections and Services}
%Every qualified student is welcome in my classroom.  As needed, I am happy to
%work with you, disability services, veterans' services, rural student services,
%etc to find reasonable accomodations. Students at this university are protected
%against sexual harassment and discrimination (Title IX), and minors have
%additional protections. \textit{As required,} if I notice or am informed
%of \textit{certain types} of misconduct, then I am required to report it
%to the appropriate authorities.  For more information on your rights as a
%student and the resources available to you, please go to the following site:
%\texttt{https://cms-test.alaska.edu/handbook/}.



\subheading{Incomplete Grade} 
Incomplete (I) will only be given in
  DMS courses in cases where
  the student has completed the majority (normally all but the last
  three weeks) of a course with a grade of C or better, but for
  personal reasons beyond his/her control has been unable to complete
  the course during the regular term. Negligence or indifference are
  not acceptable reasons for the granting of an incomplete
  grade. 

\subheading{Late Withdrawals} 
A withdrawal after the deadline
  (currently 9 weeks into the semester) from a DMS course will
  normally be granted only in cases where the student is performing
  satisfactorily (i.e., C or better) in a course, but has exceptional
  reasons, beyond his/her control, for being unable to complete the
  course. These exceptional reasons should be detailed in writing to
  the instructor, department head and dean.

\subheading{No Early Final Examinations}
 Normally, a student will not be
  allowed to take a final exam early. Exceptions can be made by
  individual instructors, but should only be allowed in exceptional
  circumstances and in a manner which doesn't endanger the security of
  the exam.

\pagebreak

 %\begin{center} \textsc{Syllabus Addendum} \end{center}
 
% \noindent{\bf COVID-19 statement:} Students should keep up-to-date on the university's policies, practices, and mandates related to COVID-19 by regularly checking this website: \url{https://sites.google.com/alaska.edu/coronavirus/uaf?authuser=0}

%Further, students are expected to adhere to the university's policies, practices, and mandates and are subject to disciplinary actions if they do not comply.


\begin{center}
\textbf{\large{Official UAF Syllabus Addendum}}
\end{center}


\noindent{\bf Student protections statement:} The university respects and upholds the principles of due process and a fair and equitable process as specified in the Board of Regents' Policy 09.02 Student Rights and Responsibilities. For more information regarding the rights and responsibilities of students, refer to the Office of Rights, Compliance and Accountability website. You are encouraged to read the Board of Regents' policy carefully to fully understand your responsibilities to our community.

We strive to create a safe and respectful environment for all members of our community. If you have questions about expectations of you as a student or believe your rights are being violated, we encourage you to reach out to the  Office of Rights, Compliance and Accountability for help. UAF reserves the right to suspend, expel or take other necessary and appropriate action in cases where a student is unable or unwilling to uphold community standards and campus safety.

For more information on your rights as a student and the resources available to you to resolve problems, please go to the following site:\\ \url{https://catalog.uaf.edu/academics-regulations/students-rights-responsibilities/}

\noindent{\bf Disability services statement:} I will work with the Office of Disability Services to provide reasonable accommodation to students with disabilities.

\noindent{\bf ASUAF advocacy statement:} The Associated Students of the University of Alaska Fairbanks, the student government of UAF, offers advocacy services to students who feel they are facing issues with staff, faculty, and/or other students specifically if these issues are hindering the ability of the student to succeed in their academics or go about their lives at the university. Students who wish to utilize these services can contact the Student Advocacy Director by visiting the ASUAF office or emailing asuaf.office@alaska.edu. 



\noindent{\bf Student Academic Support:}
\begin{itemize}
\setlength\itemsep{0em}
        \item Communication Center (907-474-7007, \mailto{uaf-commcenter@alaska.edu}, Student Success Center, 6th Floor Room 677 Rasmuson Library)
        \item Writing Center (907-474-5314, \mailto{uaf-writing-center@alaska.edu}, Student Success Center, 6th Floor Room 677 Rasmuson Library)
\item UAF Math Services (907-474-7332, \mailto{uaf-traccloud@alaska.edu})


\begin{itemize}
\item Drop-in tutoring, Student Success Center, 6th Floor Room 672 Rasmuson Library

\item 1:1 tutoring (by appointment only), 6th Floor Room 677 Rasmuson Library

\item Online tutoring (by appointment only) available

https://www.uaf.edu/dms/mathlab/, available at the Student Success Center
\end{itemize}

\item Developmental Math Lab, Gruening 406
\item The Debbie Moses Learning Center at CTC (907-455-2860, 604 Barnette St, Room 102,\\ \url{https://www.ctc.uaf.edu/student-services/student-success-center/})
\item For more information and resources, please see the Academic Advising Resource List (\url{https://www.uaf.edu/advising/students/index.php})
\end{itemize}

\noindent{\bf Student Resources:}
\begin{itemize}
\setlength\itemsep{0em}
\item Disability Services (907-474-5655, \mailto{uaf-disability-services@alaska.edu}, 110 Eielson Building)
\item Student Health \& Counseling [free counseling sessions available] (907-474-7043, \url{https://www.uaf.edu/chc/appointments.php}, Whitaker Building, Room 206, Health, Safety \& Security Bldg --- same building as Fire and Police)
\item Office of Rights, Compliance and Accountability (907-474-7300, \mailto{uaf-orca@alaska.edu}, 3rd Floor, Constitution Hall)
\item Associated Students of the University of Alaska Fairbanks (ASUAF) or ASUAF Student Government (907-474-7355, \mailto{asuaf.office@alaska.edu}{asuaf.office@alaska.edu}, Wood Center 119)
\end{itemize}

\noindent{\bf Nondiscrimination statement:}
Nondiscrimination statement: The University of Alaska is an equal opportunity/equal access employer, educational institution and provider. The University of Alaska does not discriminate on the basis of race, religion, color, national origin, citizenship, age, sex, physical or mental disability, status as a protected veteran, marital status, changes in marital status, pregnancy, childbirth or related medical conditions, parenthood, sexual orientation, gender identity, political affiliation or belief, genetic information, or other legally protected status. The University's commitment to nondiscrimination, including against sex discrimination, applies to students, employees, and applicants for admission and employment. Contact information, applicable laws, and complaint procedures are included on UA's statement of nondiscrimination available at \url{www.alaska.edu/nondiscrimination}.

\begin{tabular}{l}
UAF Office of Rights, Compliance and Accountability\\
1692 Tok Lane\\
3rd floor, Constitution Hall, Fairbanks, AK 99775\\
907-474-7300\\
\url{uaf-orca@alaska.edu}
\end{tabular}


 \scriptsize syllabus version: \today \normalsize

\end{document}

